% SPDX-License-Identifier: MIT
\chapter{Willkommen beim SASD UI Toolkit}\label{chap:intro}
\ChapterSummary{Sie lernen die Zielsetzung, die Architektur und den aktuellen
Reifegrad des Toolkits kennen. Am Ende können Sie einschätzen, für welche
Projekte es heute geeignet ist, welche Teile Sie zuerst kennenlernen sollten
und was bewusst noch nicht versprochen wird.}

\section{Wozu dieses Handbuch dient}
Eine Bibliothek wird nicht schon dadurch nützlich, dass sie viele Klassen
besitzt. Nützlich wird sie, wenn Entwickler die richtigen Bausteine finden,
das Zusammenspiel verstehen, Fehler vorhersehen und eine Oberfläche
überschaubar erweitern können. Dieses Anwenderhandbuch richtet sich an
C++-Entwickler, die \emph{Anwendungen mit} dem SASD UI Toolkit bauen wollen.
Es ist weder eine interne Entwickler-Spezifikation noch der Ersatz für die
API-Deklarationen in den Header-Dateien.

Das Buch ist wie ein klassisches Software-Handbuch aufgebaut: Zuerst
Orientierung und Installation, dann ein geführtes Praxisbeispiel, danach
die einzelnen Bedien- und Programmierkonzepte und zum Schluss
Problembehandlung und ein kompakter Nachschlageteil. Jedes Kapitel beginnt
mit einer Zusammenfassung; konkrete Quellpfade erleichtern die Vertiefung.

\section{Die Idee: eine Semantik, mehrere Darstellungen}
Ein \term{Button} bedeutet unabhängig von seinem Ausgabegerät:
Er ist sichtbar oder verborgen, aktiviert oder deaktiviert, kann den
Fokus besitzen und eine Aktion auslösen. Wie seine Beschriftung,
seine Umrandung oder sein gedrückter Zustand aussehen, bestimmt dagegen
die Darstellungsschicht.

\begin{center}
\begin{tikzpicture}[every node/.style={draw,rounded corners=2pt,align=center,
 font=\sffamily\small,minimum height=10mm},node distance=7mm]
\node[fill=SASDPale,text width=10cm] (app) {Anwendungscode: Widgets, Datenmodelle, Commands};
\node[below=of app,fill=gray!7,text width=10cm] (core) {SASD::UI Core: Layout, Fokus, Ereignisse, Lebensdauer};
\node[below left=11mm and -16mm of core,fill=gray!7,text width=3.8cm] (term) {Terminal\\ANSI/VT};
\node[below right=11mm and -16mm of core,fill=gray!7,text width=3.8cm] (rend) {Rendered\\SDL3 (experimentell)};
\draw[-stealth,thick,SASDBlue] (app)--(core);
\draw[-stealth,thick,SASDBlue] (core)--(term);
\draw[-stealth,thick,SASDBlue] (core)--(rend);
\end{tikzpicture}
\end{center}

Der Core ist \term{backendneutral}; er kennt keine Win32-Handles,
SDL-Fenster oder terminalinternen Steuersequenzen. Das Terminal
ist eine vollwertige Zielplattform. Die gerenderte Darstellung besitzt
eine eigene, auch ohne SDL verwendbare Befehlsschicht.
Native Desktop-Peers für Win32, GTK oder AppKit sind langfristige Ziele,
keine heute austauschbaren Produktiv-Backends.

\section{Was gegenwärtig vorhanden ist}
Die folgenden Aussagen beziehen sich auf den Quellstand vom
12.\ Oktober 2026; für neuere Entwicklungen sind README, Tests und
Header massgeblich.

\noindent\begin{tabularx}{\linewidth}{@{}p{32mm}X@{}}
\toprule
Bereich & Einordnung\\\midrule
Core & C++20, Komponenten- und Widget-Baum, Ereignisse, Fokus,
Layout, Commands, Shortcuts, Menumodelle, mehrere Form- und Daten-Widgets.\\
Terminal & Benutzerseitig erprobtes ANSI/VT-Backend, TerminalSession
mit Wiederherstellung, Demo und Terminaltests.\\
Rendered & Backendneutrale DisplayList, Messung und Rendering-Schnittstellen;
eigene Tests.\\
SDL3 & Optionaler experimenteller Adapter mit Fenster- und
Headless-Pfad; \emph{build-tree-only}.\\
Native Desktop & Zukünftig; kein fertiger gleichwertiger Win32-/GTK-/AppKit-Peer.\\
Stabilität & Vor Version 1.0; Quellkompatibilitaet kann sich bei
begruendeten Architekturverbesserungen ändern.\\
\bottomrule
\end{tabularx}

\section{Wichtige Begriffe}
\term{Widget} ist ein sichtbares, semantisches Bedienelement.
\term{Container} besitzt \emph{Components} und ordnet sichtbare
\emph{Widgets} als Kinder ein. \term{Backend} verbindet die
Anwendungsereignisse mit einer Zielumgebung. \term{Model} liefert
Daten, \term{View} stellt einen Ausschnitt dar. \term{Command} ist eine
nichtvisuelle, wiederverwendbare Aktion. \term{Presentation Sink}
erzeugt eine Darstellung aus dem aktuellen Zustand.

\section{Lesepfade}
Für einen schnellen Einstieg lesen Sie die Kapitel
\ref{chap:install} und \ref{chap:first}. Wer eigene Formulare bauen will,
vertieft Komponenten, Layout und Events in den Kapiteln
\ref{chap:widgets} und \ref{chap:layout}. Für administrative Anwendungen
mit Listen und Hierarchien ist Kapitel \ref{chap:models} besonders wichtig.
Kapitel \ref{chap:backends} klärt die Grenzen der Ausgabegeräte.

\PracticalTip{Beginnen Sie mit der mitgelieferten Terminal-Demo und
verändern Sie jeweils nur eine Sache: zuerst Text, dann ein Widget,
dann die Fokusfolge, zuletzt Events oder Menüs.}
